\documentclass[a4paper, twoside]{report}

%% Language and font encodings
\usepackage[english]{babel}
\usepackage[utf8x]{inputenc}
\usepackage[T1]{fontenc}

%% Sets page size and margins
\usepackage[a4paper,top=3cm,bottom=2cm,left=3cm,right=3cm,marginparwidth=1.75cm]{geometry}

%% Useful packages
\usepackage{amsmath}
\usepackage{graphicx}
\usepackage{xcolor}
\usepackage[colorinlistoftodos]{todonotes}
\usepackage[colorlinks=true, allcolors=blue]{hyperref}

\newcommand{\rcite}[1]{\textcolor{red}{\cite{#1}}}

\title{Implementation of Foundational Verification for Probabilistic Programs in Continuous Domain}
\author{Edwin Fernando}
% Update supervisor and other title stuff in title/title.tex

\begin{document}
\input{title/title.tex}

\begin{abstract}
Probabilistic Programming offers a principled way to decouple inference algorithms from the
statistical models a programmer/statistician wishes to create. There are several languages
(anglican hakaru, \rcite) used predominantly by statisticians. But there are much wider applications
in computer science, where arguments of correctness are probabilistic in nature. Examples
include Zero Knowledge Proof Systems (Blockchain), Differential Privacy and Machine Learning. We can reason 
expressively and formally verify these arguments if the implementation was carried out in a probabilistic
programming language/DSL, using Probabilistic Separation Logics (PSL).
This project is meant as a proof of concept, aiming to be the first implementation of a PSL.
Specifically in the Lean proof assistant, we implement the lilac separation logic
using Iris-lean (ported from rocq) to provide ergonomic tactical support. We choose
Lean for the significant body of measure-theoretic results present in Mathlib, required
to establish soundness of the logic.

Given time I would also like to pursue formalisation of BaSL or a
 theory for symbolic execution using lilac or BaSL for probabilistic programs.


on top of the Iris framework, in Lean.
offer an intuitive way to reason about
probabilistic programs at a suitable level of abstraction. 
\end{abstract}

% \renewcommand{\abstractname}{Acknowledgements}
% \begin{abstract}
% Thanks mum!
% \end{abstract}

\tableofcontents
% \listoffigures
% \listoftables

\chapter{Introduction}
Probabilistic programming languages (Hakaru, Stan, Gen, etc \cite{hakaru, stan, gen}) offer a principled way of 
creating statistical models complex probability distributions
over multiple random variables, as code. The programmer (statistician) can
then interact with the model, evaluating (sampling from) the generated distribution.
This decouples the specification of a model from the inference algorithm used for sampling.
Probabilistic programming is in fact also pervasive in computer science, though generally
not cleanly abstracted into a DSL. Examples include Zero Knowledge Proof Systems (Blockchain) \cite{arklib},
Differential Privacy and Machine Learning (TensorFlow Probability \cite{tfp}). Formal reasoning about programs in
these critical applications may be highly beneficial. 

The aim of this project is to extend the reach of foundational
verification in this direction. The main contributions are:
\begin{enumerate}
\item \textbf{C1}: Implementing (formalising) Lilac \cite{lilac} and BaSL \cite{basl} separation logics in the Lean theorem prover,
with "tactical support", using Iris \cite{iris}.
\item \textbf{C2}: (if time permits) Develop a theory of compositional symbolic execution (CSE) for probabilistic programs and implement using Soteira \cite{soteira}.
CSE can automate program verification.
\end{enumerate}

% {\color{red} Here on is written hypothetically if I finish ALL of my goals}

\textbf{C1} is one of the first (if not the first) foundationally verified (in Lean) implementation of a probabilistic separation logic (PSL).
There is in fact a simultaneous project formalising the Bluebell PSL \cite{bluebell-impl} for the purpose of verifying the core of a Zero Knowledge Proof
system \cite{arklib}.

\textbf{C2} would be a novel theoretical contribution. The main
innovation here, would be either
The separation logic chosen
for this may be Lilac or BaSL, but also possibly Bluebell \cite{bluebell}, to directly apply the work in an industrial formalisation project. However
Bluebell supports both unary

with various subsystems 
But there are much wider applications
in computer science, where arguments of correctness are probabilistic in nature. Examples
expressively and formally verify these arguments if the implementation was carried out in a probabilistic
programming language/DSL, using Probabilistic Separation Logics (PSL).
This project is meant as a proof of concept, aiming to be the first implementation of a PSL.
Specifically in the Lean proof assistant, we implement the lilac separation logic
using Iris-lean (ported from rocq) to provide ergonomic tactical support. We choose
Lean for the significant body of measure-theoretic results present in Mathlib, required
to establish soundness of the logic.

Given time I would also like to pursue formalisation of BaSL or a
 theory for symbolic execution using lilac or BaSL for probabilistic programs.


on top of the Iris framework, in Lean.
offer an intuitive way to reason about
probabilistic programs at a suitable level of abstraction. 
\section{Objectives}
\section{Challenges}
\section{Contributions}
\chapter{Background}

\section{Lean}

A lot of background in Lean was required to undertake this project.
The Lean development exercises the full expressivity of inductive types in
Lean, for example and a thorough understanding of propositional vs definitional
equality which is related to Proof irrelevance [Theorem proving for Lean4] and how that's implemented
in Lean's type theory by special treatment of the lowest Universe [Mario's thesis]

Various representations of formal languages particularly the problem of how
to represent variables in contexts, with strong guarantees from the types System.
[Cite chlipala, HOAS, and lean-lang.org]

\section{Probabilistic Program Semantics}
This section is entirely a summary of the denotational semantics given to Bayesian Probabilistic programs
given by Staton. The probabilistic programs we aim to provide verification tools for in this work
do not go beyond the generality of the semantics presented by Staton. A brief overview of key ideas
is provided below (but it is better to read the )
\begin{itemize}
\item A clear treatment of Staton's commutative semantics.
\item The distinction and intuition behind the difference
between s-finite and sigma-finite measures. Perhaps
even why they're used.
\end{itemize}


\section{Separation Logic}

O'Hearn's original paper. PSL, Lilac, BaSL

It is unclear to me why an axiomatic semantics would need to be provided in terms of measurable maps
from the Hilbert cube. I wonder whether measure theory is not abstract enough and perhaps synthetic
measure theory may hold some intuitions/inspiration or simplifications 
\section{Measure Theory}

\subsection{Pi-lambda theorem}

\subsection{Disintegration theorem}

\section{Symbolic Execution tools}

\section{Iris and MoSeL}


\chapter{PROJECT X}

I want to note that I have not before seen the interpretation of
a typed calculus, taking both the type and typing context as parameters
to give the denotation of a term. This definition of terms with they're typing
context, actually massively simplifies the lean development
previous iterations, involved a type inference algorithm, which could fail
and the denotation function was partial.

Possible goals of this project:
\begin{itemize}
\item Formalisation of Lilac with tactical support
\item Lean interpreter for the language which is sound w.r.t to the
denotational semantics
\item The first case of Compositional Symbolic execution for probabilistic separation logic
  \subitem CSE for the Bluebell separation logic with possible immediate application in industry
  for verifying correctness of Zero Knowledge Proof systems
  This could perhaps be a tactic, which actually invokes the symbolic execution
  to elaborate into a sequence of proof rules?

  This is possibly not so hard
  \subitem CSE for Lilac or BaSL. this is perhaps out of scope of project. But it
  connects to the first part of the project formalising Lilac/BaSL which is novel
  for being the first mechanisation of a separation logic for continuous probabilistic
  programs.
\end{itemize}

terms in
context
% \input{evaluation/evaluation.tex}
% \input{conclusion/conclusion.tex}
% \appendix
\chapter{Declarations}
\section{Use of Generative AI}
Generative AI was not used in the writing of this report. It was however used to speed up theorem proving in Lean and a discussion of my experience is provided in Sec \ref{sec:wp-meas}. The only
use of Generative AI was through \href{https://aristotle.harmonic.fun/}{Aristotle}, a tool available publicly available online, free of cost. 

\section{Ethical Considerations}
This project pushes the boundary of for formal correctness guarantees of a certain class of programs called probabilistic programs. The applications are strictly in increasing the safety of software systems,
so there is no notable ethical considerations.

\section{Sustainability}
The most computationally resource intensive tool used in this project is expected to be \href{https://aristotle.harmonic.fun/}{Aristotle}, a Generative AI based tool for writing Lean code.
The energy usage of this tool is not publicly made available, however it is completely free to the public, which suggests that its energy usage cannot be too high. In line with the goals of this
project for extending the Lean development into more complex logics, Aristotle was only used ``trivial constructions'' or ``obvious sub-lemmas'', since it produces convoluted proofs which are not line with our
requirements for more significant tasks. A detailed discussion can be found in Sec \ref{sec:wp-meas}. But the takeaway here is that the principle followed, should curtail compute-intensive situations
where the system is unable to solve the tasks and goes in circles till timeout. 
Other than that, and the compilation of Lean and Latex files on a local PC, there are no other sustainability considerations.

\section{Availability of Materials}
This report is submitted with a GitHub repository containing the Lean artefact, which is made publicly available: \url{https://github.com/edwin1729/bppl-logic/tree/submission}.

\bibliographystyle{alpha}
\bibliography{bibs/sample}

\end{document}